\scanpage{73}
% archive foliation: 72
recollement de ces morphismes est immédiat.
Par suite, on a défini $\underline W(X)\xrightarrow{\varphi}\mathfrak X$.
Je dis que c'est une solution du Pb universel, i.e. que pour tt $T$,
\[
\mathrm{Hom}_S(T,X)\longrightarrow
\mathrm{Hom}_{\underline W(S)}(\underline W(T),\mathfrak X)
\]
est \uline{bijectif}. Soit \deleted{$u:T\to X$},
\[
\underline W(T)\xrightarrow{v}\mathfrak X
\]
je dis qu'il peut se factoriser de \uline{façon unique} en
\diagram{\begin{tikzcd}
\underline W(T) \arrow[r,"\underline W(u)"] &
\underline W(X) \arrow[r,"\varphi"] & \mathfrak X
\end{tikzcd}}
En effet, \deleted{\unsure{considérons} \unsure{les} \unsure{images}}
\deleted{\unsure{supposons} \unsure{qu'il} \unsure{se} \unsure{peut factoriser}, \unsure{et} \unsure{soit} $T=\bigcup$}
\deleted{$\underline W(u)($} \deleted{\unsure{inverse}}
utilisons la condition $(\beta)$ : $T=\bigcup T_i$,
avec $v(\underline W(T_i))\subset\mathfrak X_i$. Donc
$v|\underline W(T_i)\unsure{)}$ est défini par $u_i:T_i\to X_i$,
par factorisation
\diagram{\begin{tikzcd}
\underline W(T_i) \arrow[r,"\underline W(u_i)"] &
\underline W(X_i) \arrow[r,"\varphi_i"] & X_i
\end{tikzcd}}
On vérifie aussitôt que les composés
\diagram{\begin{tikzcd}
T_i \arrow[r] & X_i \arrow[r] & X
\end{tikzcd}}
\deleted{\diagram{\begin{tikzcd}
\underline W(T_i) \arrow[r] & \underline W(X_i) \arrow[r] & \underline W(X)
\end{tikzcd}}}
